\documentclass[conference]{IEEEtran}
\usepackage{amsmath}
\usepackage{booktabs}
\usepackage{graphicx}
\usepackage{url}

\title{A Reproducible Fraud-Scoring Demonstrator with Explicit Audit-Write Provenance}
\author{\IEEEauthorblockN{Author names withheld for draft review}
\IEEEauthorblockA{Department of Computer Science and Engineering\\Tula's Institute, Dehradun}}

\begin{document}
\maketitle

\begin{abstract}
This paper documents an academic prototype for rare-event credit-card fraud scoring and audit-record demonstration. The application combines a multilayer perceptron (MLP), a Random Forest, and a Flask interface that can request a local Ganache smart-contract write for a flagged transaction. The original training artifacts were produced with preprocessing fitted before the train/test split, so their historical metrics are not used as evidence here. We specify a revised stratified train/validation/test protocol with train-only preprocessing, validation-based threshold selection, and a held-out comparison using precision, recall, F1, PR-AUC, and ROC-AUC. The dashboard's repeated feed is generated by perturbing stored public dataset examples; it is not a live payment stream. Ganache receipts and volatile local simulation records are reported as distinct outcomes. Revised metrics and local scoring measurements are reported from actual runs on the documented setup.
\end{abstract}

\begin{IEEEkeywords}
credit-card fraud detection, class imbalance, model evaluation, audit logging, blockchain prototype
\end{IEEEkeywords}

\section{Introduction}
Fraud detection is a rare-event classification problem in which accuracy alone can obscure the number of missed fraud cases and the cost of false alerts. Evaluation choices such as the operating threshold and the class distribution used for training affect how scores should be interpreted. Prior work has examined realistic operational modeling, class imbalance, and probability calibration in credit-card fraud detection \cite{dalpozzolo2018,dalpozzolo2015}.

This work describes GuardChain AI, a local software prototype that scores one transaction feature vector per request and presents a fraud audit workflow. Its scope is deliberately limited: the dataset is historical and anonymized, the dashboard simulator perturbs stored examples, and the blockchain network is local Ganache. The prototype is not evaluated as a production payment service.

\subsection{Contributions}
The engineering and evaluation contributions are:
\begin{itemize}
  \item A reproducible 65/15/20 stratified split protocol with train-only scaling and validation-only threshold selection for the MLP and Random Forest.
  \item A like-for-like held-out comparison reporting precision, recall, F1, PR-AUC, ROC-AUC, confusion matrices, and class support.
  \item An application response that distinguishes confirmed Ganache writes from process-local simulation records and records model/scoring timing separately from audit-write timing.
  \item A self-contained dashboard path that demonstrates analyzer and sample-replay flows when external transaction infrastructure is unavailable.
\end{itemize}

These are prototype and evaluation contributions. They do not establish a new classifier, a production-grade ledger, or a novel blockchain consensus design.

\section{Related Work}
Dal Pozzolo et al. describe operational considerations in card-fraud detection, including class imbalance, concept drift, delayed labels, and realistic evaluation \cite{dalpozzolo2018}. Their work motivates reporting measures beyond global accuracy and evaluating the system under realistic data assumptions. Dal Pozzolo et al. also analyze how undersampling can bias posterior probabilities and why decision thresholds must be considered with the sampling scheme \cite{dalpozzolo2015}. The present prototype uses class-weighted loss and validation-based threshold choice; it does not claim probability calibration or robustness to concept drift.

The supplied project report discusses blockchain-backed audit records, but the inspected implementation uses a single local Ganache node and stores sender/receiver strings and amounts in contract state. This is an implementation observation, not evidence of multi-party trust, privacy, or production immutability.

\section{Methodology}
\subsection{Dataset and Inputs}
The project uses the public European cardholder transaction dataset distributed by the Machine Learning Group at Universit\'e Libre de Bruxelles (ULB) as \texttt{creditcard.csv} \cite{ulbdataset}. It contains 284,807 rows, 30 input features (\texttt{Time}, \texttt{Amount}, and \texttt{V1}--\texttt{V28}), and binary label \texttt{Class}. The local data file used for this run is identified by its SHA-256 checksum in the experiment artifacts; results should be reproduced only with a matching file.

\subsection{Models and Split Protocol}
The MLP comprises three fully connected hidden layers with Batch Normalization, ReLU, and Dropout, followed by one logit. The loss is binary cross-entropy with logits; the positive weight is calculated from training rows. The Random Forest is a class-weighted tabular baseline. Its number of trees is selected between 50 and 100 by validation PR-AUC; the selected run used 100 trees (validation PR-AUC 0.8245 versus 0.8120 for 50 trees). A fixed seed is used to create stratified train, validation, and test partitions with proportions 0.65, 0.15, and 0.20. Scalers are fitted on training rows and applied unchanged to validation and test rows. The decision threshold is selected by maximizing validation F1 and then held fixed for test evaluation.

\subsection{System and Audit Flow}
The browser sends numeric feature values to Flask. The server validates the feature set, runs the selected model, and returns the probability, threshold, and timing fields. For a fraud verdict, Flask invokes the Web3 writer. If Ganache is reachable and the contract call succeeds, the response is marked as Ganache; otherwise, the record is labelled as a volatile local simulation. A synthetic simulation identifier is not presented as an Ethereum transaction hash.

\section{Experimental Setup}
Experiments are implemented in \texttt{train\_dl\_model.py}, \texttt{train\_model.py}, \texttt{ablation.py}, and \texttt{evaluate.py}. The scripts save model artifacts, split metadata, probabilities, metrics, and plots under \texttt{artifacts/phase3/} and \texttt{results/phase3/}. The run used Windows 11, Python 3.12.14, 12 logical CPU cores, 7.7 GiB reported system memory, and CPU-only PyTorch 2.14.1. Package versions include NumPy 2.5.4, pandas 3.1.0, scikit-learn 1.9.1, Flask 3.1.3, and Web3.py 7.16.0. The dataset SHA-256 is \texttt{76274b691b16a6c49d3f159c883398e03ccd6d1ee12d9d8ee38f4b4b98551a89}.

For local scoring performance, the Flask test client measures requests with audit writes stubbed. These values exclude network-client and Ganache confirmation time. Ganache transaction latency must be reported separately from model inference and Flask scoring.

\section{Results}
Metrics below were produced by the revised run. The operating threshold was selected on validation data by maximizing F1 and then held fixed for the test partition.

\begin{table*}[htbp]
\caption{Held-out comparison on 56,962 transactions (98 fraud cases).}
\centering
\begin{tabular}{lccccccc}
\toprule
Model & Threshold & Accuracy & Precision & Recall & F1 & PR-AUC & ROC-AUC \\
\midrule
MLP & 0.9798 & 0.9992 & 0.7217 & 0.8469 & 0.7793 & 0.6745 & 0.9746 \\
Random Forest & 0.3600 & 0.9996 & 0.9419 & 0.8265 & 0.8804 & 0.8609 & 0.9631 \\
\bottomrule
\end{tabular}
\end{table*}

At the selected thresholds, the MLP confusion matrix is 56,832 true negatives, 32 false positives, 15 false negatives, and 83 true positives. The RF matrix is 56,859 true negatives, 5 false positives, 17 false negatives, and 81 true positives.

\begin{table}[htbp]
\caption{MLP class-weight ablation on the held-out partition.}
\centering
\begin{tabular}{lcccc}
\toprule
Loss & Threshold & Precision & Recall & F1 \\
\midrule
Weighted BCE & 0.9798 & 0.7217 & 0.8469 & 0.7793 \\
Unweighted BCE & 0.5175 & 0.8163 & 0.8163 & 0.8163 \\
\bottomrule
\end{tabular}
\end{table}

In this single split, weighted loss raises recall by 3.06 percentage points but lowers precision and F1; PR-AUC is 0.6745 for weighted and 0.7445 for unweighted BCE. This does not establish a general advantage for either loss.

The local Flask test-client benchmark discarded ten warmup calls per model and measured 490 requests per model with audit writes stubbed. MLP scoring p50/p95 was 1.82/3.01 ms at 389.6 requests/s; RF scoring p50/p95 was 35.76/49.19 ms at 26.2 requests/s. These measurements exclude network transport and Ganache confirmation. Ganache confirmation latency is not reported.

\section{Discussion}
The software path demonstrates integration between a classifier and an audit-record interface, but the simulator does not establish live transaction ingestion. A local Ganache chain demonstrates contract calls, not decentralized multi-bank governance. Persisting payment identifiers or transaction details on a public or consortium chain can create privacy and retention concerns. The revised evaluation is designed to expose false-alert workload and to prevent test-set preprocessing leakage.

\section{Limitations}
The benchmark dataset is historical and contains anonymized PCA features; it does not represent a contemporary bank's production feature stream. There is no external payment gateway, online concept-drift evaluation, calibrated risk model, user authentication, role-based case management, durable simulation queue, or multi-node blockchain deployment. No production security, compliance, or real-time service-level claim is made.

\section{Conclusion}
GuardChain AI is an educational fraud-scoring and audit-workflow prototype. On this held-out split, the Random Forest has higher precision and PR-AUC, while the MLP has higher recall; weighted loss modestly shifts that MLP operating point toward recall. The local scoring benchmark is below the 500 ms target for both models, but does not include a real client network or chain confirmation. The simulator remains sample replay, and these results do not establish production behavior.

\begin{thebibliography}{00}
\bibitem{dalpozzolo2018} A. Dal Pozzolo, G. Boracchi, O. Caelen, C. Alippi, and G. Bontempi, ``Credit Card Fraud Detection: A Realistic Modeling and a Novel Learning Strategy,'' \emph{IEEE Transactions on Neural Networks and Learning Systems}, vol. 29, no. 8, pp. 3784--3797, 2018, doi: 10.1109/TNNLS.2017.2736643.
\bibitem{dalpozzolo2015} A. Dal Pozzolo, O. Caelen, R. A. Johnson, and G. Bontempi, ``Calibrating Probability with Undersampling for Unbalanced Classification,'' in \emph{Proc. IEEE Symposium Series on Computational Intelligence}, 2015, pp. 159--166, doi: 10.1109/SSCI.2015.33.
\bibitem{ulbdataset} Machine Learning Group, Universit\'e Libre de Bruxelles, ``Credit Card Fraud Detection Dataset,'' Kaggle. [Online]. Available: \url{https://www.kaggle.com/datasets/mlg-ulb/creditcardfraud}. Accessed: Oct. 23, 2026.
\end{thebibliography}

\end{document}
